\chapter{Optional Libraries to be Linked into FHI-aims}
\label{appendix_external_packages_and_variants}

A core principle for FHI-aims is that it should be possible to build
all essential functionality with only a minimal set of dependencies,
typically a Fortran compiler, BLAS and LAPACK libraries, an MPI
library, and the scaLAPACK library. These packages are essential for
performance and usually available on every relevant supercomputer
architecture. 

We note that while it is often easy to fulfil other,
far more complex dependencies on standard Linux type systems with
household tools, this is not always true for highly specialized
supercomputers. Yet, that computer may be precisely one the
multi-million dollar machine which would help solve an advanced
scientific problem. It is essential that a mere human be able to
compile FHI-aims with the appropriate performance on such a machine. 

That said, it may be beneficial to connect other libraries to FHI-aims
for special purposes, including libraries which require
cross-compilation with a C compiler. This section is intended to be
developed into a list of such libraries. Please also see Section
\ref{cross-compile-c} for some basic information.

\section{Adding Optional Libraries into FHI-aims: Stubs}

The optional presence or absence of certain libraries from FHI-aims at
compile time means that the main FHI-aims code must be able to deal
with the possible absence of certain subroutine or function calls to
those libraries at compile time.

In many codes, this is handled by adding preprocessor statement. In
FHI-aims, preprocessor flags are \textbf{NOT} the way forward. The
simple reason is that a single preprocessor flag may remain
understandable for an ordinary user; the proliferation of ten or more
preprocessor flags will create a serious obstacle for others to obtain
a reasonably compiled code. Code that introduces preprocessor
statements into FHI-aims will be removed from the code base.  

In FHI-aims, the recommended way to optionally add or circumvent a
given library is by creating a ``stub'' file for those library calls
-- essentially, empty subroutines that inform a user that they should
not have ended up here with a given version of the code (i.e., without
linking to a certain external library). At compile time, an
appropriate flag in the Makefile and Makefile.backend should be
created that switches between the real external library (if available)
or the ``stub'' file. Please follow the example of ``Spglib'' below in
the Makefile and in Makefile.backend for more information. The
principle is simple and easy to apply for any other optional library.

\section{Spglib}

Spglib is a library for finding and handling crystal
symmetries, written by Atsushi Togo
(\textbf{http://spglib.sourceforge.net}). FHI-aims can be interfaced
to the spglib.  

So far, FHI-aims only supports the determination
of the crystal symmetry for a given 
(periodic) geometry file, but does not make use of this information
otherwise. 

Prerequisites: \\

spglib is written in C and therefore needs a C compiler with
appropriate options in addition to the usual Fortran compiler with
which FHI-aims is built. Please enable a C compiler and add the
appropriate compiler flags to the \texttt{Makefile} or
\texttt{make.sys} file as described in Section \ref{cross-compile-c}. 

Specifically for spglib, include in \texttt{Makefile} or \texttt{make.sys}:\\
  
\begin{verbatim}
    USE_SPGLIB = yes
\end{verbatim}



\keydefinition{use\_symmetry\_analysis}{control.in}
{
  \noindent
  Usage: \keyword{use\_symmetry\_analysis} \option{.true. / .false.} \\[1.0ex]  
  Purpose: This flag activates the interface to the spglib to obtain 
  symmetry information for the provided geometry file once.
  The symmetry information is directly written after the geometry printout
  in the FHI-aims standard output. spglib needs to be compiled with FHI-aims 
  to get the output. \\[1.0ex]
  Default: \option{.true.} \\
}
\keydefinition{sym\_precision}{control.in}
{
  \noindent
  Usage: \keyword{sym\_precision} \texttt{value} \\[1.0ex]  
  Purpose: This value determines the accuracy for identifying symmetric atoms. 
  The lower the value, the more strictly the atomic positions must overlap.\\[1.0ex] 
  Default: $10^{-5}$
}
